\begin{tikzpicture}[x=1mm,y=1mm,font=\figurefont\footnotesize,text=NNInk,
  nn flow/.append style={draw=NNWire,rounded corners=1.5mm},
  nn operator path/.append style={draw=NNWire,rounded corners=1.5mm},
  nn operator/.append style={fill=NNHidden!18}]
  \path[nn frame] (18,-70) rectangle (97,26);
  \node[anchor=east] (memory) at (12,17) {$\bm h_{1:T}$};
  \node[anchor=east] (old) at (12,-40) {$\bm s^{(j-1)}$};
  \node[anchor=east] (input) at (12,-62) {$\bm e(y_{j-1})$};
  \node[anchor=west] (state) at (99,-24) {$\bm s^{(j)}$};
  \node[anchor=west] (output) at (99,-40) {$\bm p^{(j)}$};
  \node[nn operator,minimum width=28mm,minimum height=12mm] (score) at (46,0) {评分＋Softmax};
  \node[nn operator,draw=NNOutput,fill=NNOutput!18,minimum width=24mm,minimum height=12mm] (ctx) at (81,0) {加权聚合};
  \node[nn operator,minimum width=25mm,minimum height=12mm] (dec) at (46,-40) {解码更新};
  \node[nn operator,draw=NNOutput,fill=NNOutput!18,minimum width=18mm,minimum height=12mm] (g) at (81,-40) {输出映射};
  % 两条输入总线分别供注意力读取和目标嵌入，分支保留清晰的连接点。
  \draw[nn flow] (memory.east)--(81,17)--(ctx.north);
  \draw[nn flow] (46,17)--(score.north);
  \fill[NNWire] (46,17) circle[radius=.45mm];
  \draw[nn flow] (old.east)--(dec.west);
  \draw[nn flow] (27,-40)--(27,0)--(score.west);
  \fill[NNWire] (27,-40) circle[radius=.45mm];
  \draw[nn flow] (input.east)--(81,-62)--(g.south);
  \draw[nn flow] (46,-62)--(dec.south);
  \fill[NNWire] (46,-62) circle[radius=.45mm];
  % 权重向右交给聚合算子，上下文向下同时供解码与读出。
  \draw[nn operator path] (score.east)--node[above] {$\bm\alpha_j$} (ctx.west);
  \draw[nn operator path] (ctx.south)--(g.north);
  \draw[nn operator path] (81,-15)--node[above] {$\bm c^{(j)}$} (46,-15)--(dec.north);
  \fill[NNWire] (81,-15) circle[radius=.45mm];
  \draw[nn operator path] (dec.east)--(g.west);
  % 新状态跨过上下文支路时留出间隙，避免被误读为汇合。
  \draw[nn operator path,preaction={draw=NNPanel,line width=2.6pt,-}]
    (64,-40)--(64,-24)--(state.west);
  \fill[NNWire] (64,-40) circle[radius=.45mm];
  \draw[nn flow] (g.east)--(output.west);
\end{tikzpicture}
